\documentstyle[%


righttag%


,11pt%


,amssymb%


,russian%               remove in Eng.


,izvru%                 Set izven% in Eng.


]{amsart}





% kor@math.rsu.ru


% change \baselinestretch to fit item 8 of theorem 3 (p.7)


\newcommand{\proj}{\operatornamewithlimits{proj}}


\newcommand{\ind}{\operatornamewithlimits{ind}}


\newcommand{\co}{\operatorname{co}}


\newcommand{\la}{\langle}


\newcommand{\ra}{\rangle}


\newcommand{\tts}[1]{}





\theoremstyle{definition}


\theorembodyfont{\rm}


\newtheorem{cornonum}{\hskip\parindent Следствие}


\newtheorem{notenonum}{\hskip\parindent Замечание}


\renewcommand{\thecornonum}{}


\renewcommand{\thenotenonum}{}





%\hoffset=-15mm%                  For Eng.


\setcounter{page}{19}


\god{2005}


\nomer{9~(520)} %                        Remove in Eng.


%\nomer{Vol.\,49, No.\,9}%               Set in Eng.


%\str{\pageref{begin}--\pageref{end}}%  Set in Eng.


\udk{517.537}





\author{\it Ю.Ф.\,Коробейник}


% NB:   ^^ remove in Eng.


\title{Линейные преобразования представляющих\\ и


абсолютно-представляющих систем}


\thanks{Данная работа выполнена при


финансовой поддержке Российского фонда


фундаментальных исследований (проект \No\,02-01-00372).}





\date{}


%\pagestyle{myheadings}%         Set in Eng.


\pagestyle{plain} %              Remove in Eng.


\begin{document}


%\thispagestyle{plain}%          Set in Eng.


\maketitle


\label{begin}











\section{Основные определения. Постановка задачи}





1. Пусть $H$ --- полное отделимое локально выпуклое пространство


(ПОЛВП)


с топологией, определяемой набором преднорм $P=\{p\}$. Пусть, далее,


$X=(x_k)_{k=1}^{\infty}$ --- последовательность ненулевых элементов из $H$


такая,


что система элементов $(x_k)_{k=1}^{n}$ линейно независима в $H$ при


$n=1,2,\dotsc$ Образуем пространство числовых последовательностей


$$


A_1=A_1(X,H):= \bigg\{c=(c_k)_{k=1}^{\infty},\ c_k\in\Bbb{C}:


\ \text{ряд}\ \sum_{k=1}^{\infty}c_kx_k\ \text{сходится в}\ H\bigg \}.


$$


Положим $\forall p\in P$, $\forall c\in A_1$\;


$q_p(c):=\sup\limits_{n\geqslant 1}


p \Big(\sum\limits_{k=1}^{n}c_kx_k\Big)$ и введем в $A_1$ топологию


набором преднорм


$Q_P= \{q_p \}_{p\in P}$. Тогда $A_1$ --- ПОЛВП (см., напр., [1],


с.\,194). При этом оператор


$L:\forall c\in A_1\longrightarrow


Lc:=\sum\limits_{k=1}^{\infty}c_kx_k\in H$


является линейным непрерывным оператором из $A_1$ в $H$ (обычно он


называется {\it оператором представления}\/).





Последовательность $X$ называется {\it представляющей системой} (ПС) в


$H$, если


любой элемент $x\in H$ можно представить в виде сходящегося в $H$ ряда по


системе


$X$: $x=\sum\limits_{k=1}^{\infty}c_kx_k$, $c_k\in \Bbb{C}$,


$k\geqslant 1$.


Очевидно, $X$ --- ПС в $H$ тогда и только тогда, когда $L$ ---


эпиморфизм $A_1$ на $H$ (см., напр., [1], [2]).





В связи с тем, что разложение по $X$ любого элемента из $H$ не обязательно


единственно, имеют смысл следующие определения.





1. ПС $X$ называется {\it эффективной} (ЭПС), если для любого элемента


$x$ из $H$ имеется способ конструктивного определения коэффициентов 


какого-либо его разложения в сходящийся в $H$ ряд по системе $X$.





2. ПС $X$ называется {\it правильной} (ППС), если соответствующий оператор


представления $L$ имеет линейный непрерывный правый обратный (ЛНПО).





3. ПС $X$ называется {\it эффективно правильной} (ЭППС), если оператор


$L$ обладает эффективно определяемым ЛНПО.





Заметим, что если оператор $L$ имеет ЛНПО $T$, то $\forall x\in H$ ряд


$\sum\limits_{k=1}^{\infty}(Tx)_kx_k$ сходится в $H$ и его сумма равна $x$.


Так как топология в $A_1$ мажорирует топологию покоординатной сходимости,


то $\forall k\geqslant 1$ \, $\varphi_k(x):=(Tx)_k\in H'$, и коэффициенты


данного


разложения $x$ зависят от $x$ линейно и непрерывно. Очевидно также, что


любая ЭППС является также и ЭПС.


\medskip





2. Напомним определение одного важного подкласса ПС. Последовательность


$X=(x_k)_{k=1}^{\infty}$ ненулевых элементов $H$ называется {\it абсолютно


представляющей системой} (АПС) в ПОЛВП $H$, если любой элемент $x$ из $H$


представляется рядом


\begin{equation}


x=\sum_{k=1}^{\infty}c_kx_k, \quad c_k\in\Bbb{C}, \ \


k=1,2,\dotsc,


\end{equation}


абcолютно сходящимся в $H$. Положим


$$


A_2=A_2(X,H):= \bigg\{c=(c_k)_{k=1}^{\infty}:q_p^o(c):=


\sum_{k=1}^{\infty}|c_k|p(x_k)<\infty \ \forall p\in P \bigg\}.


$$


Тогда $A_2$ --- также ПОЛВП с набором преднорм $ \{q_p^o \}_{p\in P}$, а


оператор $L_0c:=\sum\limits_{k=1}^{\infty}c_kx_k$ линеен и непрерывен (из


$A_2$ в $H$). При этом $X$ -- АПС в $H$ тогда и только тогда, когда $L_0$ ---


эпиморфизм $A_2$ на $H$ (по поводу этих утверждений относительно $A_2$ и


$L_0$ см., напр., [1], с.\,196--197).





Так как здесь разложение в ряд (1) в общей ситуации не обладает свойством


единственности, то и в этом случае можно дать (точно так же, как для ПС, с


заменой


$L$ на $L_0$ и ПС на АПС) определения эффективной, правильной и эффективно


правильной АПС (соответственно ЭАПС, ПАПС, ЭПАПС).


\medskip





3. В данной работе исследуется 


задача, частично решенная в [3]:





{\it пусть $H_j$, $j=1,2$, --- ПОЛВП и $M$ --- линейный непрерывный оператор


из $H_1$ в $H_2$. Пусть, далее, $X$ принадлежит одному из шести


вышеперечисленных


классов ПС в $H_1$. При каких условиях последовательность


$MX:= \{Mx_k \}_{k=1}^{\infty}$ принадлежит тому же классу ПС в $H_2$}?








\section{Общие результаты}








1. Пусть $H_j$ , $j=1,2$, --- ПОЛВП с набором преднорм


$P_j=\{p_{\alpha}^j\}$,


$\alpha\in\Omega_j$; $X=(x_k)_{k=1}^{\infty}$ --- последовательность


некоторых


элементов из $H_1$, любая конечная часть которой линейно независима в $H_1$


(в результатах, относящихся к АПС, это предположение заменяется более слабым


$x_k\neq 0$ \, $\forall k\ge 1$). Всюду далее $M$ --- линейный непрерывный


оператор из $H_1$ в $H_2$. Образуем пространство последовательностей


$A_1(MX,H_2)$, где


$$


MX:=(Mx_k)_{k=1}^{\infty}.


$$








\begin{prop}


$A_1(X,H_1)\hookrightarrow A_1(MX,H_2)$.


\end{prop}








{\bf Доказательство.} Если ряд $\sum\limits_{k=1}^{\infty}d_kx_k$


сходится в $H_1$, то


в силу непрерывности и линейности оператора $M$ ряд


$\sum\limits_{k=1}^{\infty}


d_kMx_k$ сходится в $H_2$. Поэтому $A_1(X,H_1)\subseteq A_1(MX,H_2)$.


Возьмем теперь


любую преднорму $p^{(2)}\in P_2$ в $H_2$ и найдем число $b<+\infty$ и


преднорму


$p^{(1)}\in P_1$ в $H_1$ такие, что $p^{(2)}(Mv)\le bp^{(1)}(v)$ \, $\forall


v\in H_1$. Тогда $\forall d\in A_1(X,H_1)$


$$


q^1_{p^{(2)}}(d):=\sup\limits_{n\ge 1}p^{(2)}


\bigg(\sum\limits^n_{k=1}d_kMx_k\bigg)\le b\sup\limits_{n\ge


1}p^{(1)}\bigg(\sum\limits^n_{k=1}


d_kx_k\bigg)=bq_{p^{(1)}}(d).\quad\square


$$








\begin{cornonum}


Если $A_1(MX,H_2)\subseteq A_1(X,H_1)$, то $A_1(MX,H_2)=A_1(X,H_1)$.


\end{cornonum}





Аналогично доказывается








\begin{prop}


$A_2(X,H_1) \hookrightarrow A_2(MX,H_2)$.


\end{prop}





Будем говорить, что последовательность $X$ \ {\it $M$-инвариантна


относительно


пары $H_1,\ H_2$}, если $A_1(MX,H_2)=A_1(X,H_1)$, и {\it $M$-абсолютно


инвариантна


относительно той же пары}, если $A_2(MX,H_2)=A_2(X,H_1)$.








\begin{prop}


Пусть последовательность $X$ \ $M$-инвариантна $($или


$M$-абсолютно инвариантна\/$)$ относительно $(H_1,H_2)$. Пусть, далее,


$MX$ --- ПС


$($соответственно АПС\/$)$ в $H_2$. Тогда $M$ --- эпиморфизм $H_1$ на $H_2$.


\end{prop}








\begin{pf}


Предположим, что $MX$ --- АПС в $H_2$


и $X$ \ $M$-абсолютно


инвариантна относительно пары $(H_1,H_2)$. Пусть еще $y$ --- произвольный


элемент из $H_2$. Тогда $\exists(d_k)^{\infty}_{k=1}\in A_2(MX,H_2):


y=\sum\limits_{k=1}^{\infty}d_kMx_k$. Так как $A_2(MX,H_2)=A_2(X,H_1)$,


то $(d_k)_{k=1}^{\infty}\in A_2(X,H_1)$ и поэтому ряд


$\sum\limits_{k=1}^{\infty}


d_kx_k$ абсолютно сходится в $H_1$. Если $v$ --- его сумма, то $v\in H_1$


и


$$


Mv=M\bigg(\sum_{k=1}^{\infty}d_kx_k\bigg)=\sum_{k=1}^{\infty}d_kMx_k=y.


$$


Точно так же рассматривается случай, когда $MX$ --- ПС в $H_2$


и $X$ \ $M$-инвариантна относительно $(H_1,H_2)$.


\end{pf}





Аналогично доказывается








\begin{prop}


Пусть последовательность $X$ \ $M$-инвариантна


$($или $M$-абсолютно инвариантна\/$)$ относительно ($H_1,H_2$) и пусть $MX$ ---


ЭПС $($соответственно


ЭАПС\/$)$ в $H_2$. Тогда $M$ --- эпиморфизм $H_1$ на $H_2$, имеющий эффективно


определяемый правый обратный.


\end{prop}





Как отмечалось в [3], если $X$ --- ПС или ЭПС в $H_1$, а $M$ --- эпиморфизм


$H_1$ на $H_2$, то $MX$ --- ПС (соответственно ЭПС) в $H_2$. Точно так же,


если $X$ --- АПС или ЭАПС в $H_1$, а $M$ --- эпиморфизм $H_1$ на $H_2$, то $MX$


--- АПС (соответственно ЭАПС) в $H_2$. Учитывая еще предложения 3 и 4,


получаем








\begin{cor}


Пусть последовательность $X$ \ $M$-инвариантна относительно пары


$(H_1,H_2)$ и пусть $X$ --- ПС (или ЭПС) в $H_1$. Тогда равносильны такие


утверждения:


1)~$MX$ --- ПС (соответственно ЭПС) в $H_2$; 2)~$M$ --- эпиморфизм $H_1$


на $H_2$


(соответственно эпиморфизм $H_1$ на $H_2$ с эффективно определяемым


правым обратным оператором).


\end{cor}








\begin{cor}


Пусть последовательность $X$ \ $M$-абсолютно инвариантна относительно


$(H_1,H_2)$ и пусть $X$ --- АПС (или ЭАПС) в $H_1$. Тогда утверждение~2)


следствия~1 равносильно тому, что $MX$ --- АПС (или ЭАПС) в $H_2$.


\end{cor}





Назовем пару ПОЛВП $H_1$, $H_2$ {\it парой Майкла}, если любой эпиморфизм


$H_1$ на


$H_2$ обладает непрерывным (но не обязательно линейным) правым обратным


оператором


(ПОО). Как известно, если оба пространства $H_j$ являются пространствами Фреше


или оба они --- $LN^*$-пространства (определение последних см., напр., в


[4], гл.\,XI или в [5]), то согласно [6] и [5] $H_1$, $H_2$ --- пара Майкла.


Это определение позволяет дополнить предложение 3 и cледствия 1, 2.


\medskip








{\bf Предложение $\bold 3'$.} {\it Пусть выполнены исходные


предположения предложения $3$ и пусть $H_1$ и $H_2$ --- пара 


Майкла. Тогда $M$ --- эпиморфизм $H_1$ на $H_2$,


имеющий непрерывный ПОО.}


\medskip








{\bf Следствие $\bold 1'$.} Пусть последовательность


$X$ \ $M$-инвариантна


относительно пары $H_1,\ H_2$ и пусть $X$ --- ПС в $H_1$. Пусть, далее,


$H_1,\ H_2$ --- пара Майкла. Тогда $MX$ --- ПС в $H_2$ в том и только том


случае, если выполнено условие


$3'$)~$M$ --- эпиморфизм $H_1$ на $H_2$ с непрерывным ПОО.


\medskip








{\bf Следствие $\bold 2'$.} Пусть выполнены предположения


следствия~2 предложения~4 и пусть


$H_1,\ H_2$ --- пара Майкла. Тогда условие~$3'$)


следствия~$1'$ равносильно тому, что $MX$ --- АПС в $H_2$.


\medskip








2. Обратимся теперь к правильным и эффективно правильным системам.


В [3] доказано, что если $X$ --- ППС (или ЭППС) в $H_1$, а $M$ ---


эпиморфизм


$H_1$ на $H_2$, имеющий (конструктивно определяемый) линейный непрерывный


правый


обратный (ЛНПО), то $MX$ --- ППС (соответственно ЭППС) в $H_2$. Установим


теперь утверждение в известном смысле обратного характера.








\begin{prop}


Пусть последовательность $X$ \ $M$-инвариантна относительно


$(H_1,H_2)$ и $A_1(MX,H_2)$ топологически изоморфно $A_1(X,H_1)$.


Пусть, далее, $MX$ --- ППС в $H_2$. Тогда $M$ --- эпиморфизм $H_1$ на $H_2$,


имеющий ЛНПО.


\end{prop}








{\bf Доказательство.} Так как $MX$ --- ППС в $H_2$, то существует ЛНПО


$(L_1^M)_R^{-1}$ для оператора представления


$L_1^M:A_1(MX,H_2)\to H_2 \mid \ \forall d\in


A_1(MX,H_2)\longrightarrow


\sum\limits_{k=1}^{\infty}d_kMx_k\in H_2$.


При этом $\forall y\in


H_2$ \ $y=\sum\limits_{k=1}^{\infty}((L_1^M y)_R^{-1})_k Mx_k$ и этот ряд


сходится в $H_2$. Но тогда ряд


$\sum\limits_{k=1}^{\infty}((L_1^M y)_R^{-1})_k x_k$ в


силу совпадения пространств $A_1(MX,H_2)$ и $A_1(X,H_1)$ сходится в $H_1$.


Если


$v$ --- его сумма, то $Mv=y$. При этом $v=By$, где $B$ --- линейный


оператор из $H_2$ в $H_1$.





Покажем, что он непрерывен. Имеем


$$


\forall y \in H_2\ \ 


y=\sum_{k=1}^{\infty}((L_1^M y)_R^{-1})_k Mx_k,\quad


By=\sum_{k=1}^{\infty}((L_1^M y)_R^{-1})_k x_k,\quad v\in H_1.


$$


Далее $\forall p^{(1)}\in P_1$ \; $\exists b<\infty$,


$\exists p^{(2)}\in P_2$ \,$\forall


d=(d_k)_{k=1}^{\infty}\in A_1(MX,H_2)$


$$


\sup\limits_{n\ge 1}p^{(1)}\bigg(\sum_{k=1}^{n}d_kx_k\bigg)\le


b\sup\limits_{n\ge 1}p^{(2)}\bigg(\sum_{k=1}^n d_kMx_k\bigg)


$$


(в силу топологического изоморфизма $A_1(MX,H_2)$ и $A_1(X,H_1)$). Отсюда


\begin{multline*}


p^{(1)}(v)=p^{(1)}\bigg(\sum_{k=1}^{\infty}((L_1^M y)_R^{-1})_k x_k\bigg)\le\\


%


\le\sup\limits_{n\ge 1}p^{(1)}\bigg(\sum_{k=1}^{n}((L_1^M y)_R^{-1})_kx_k\bigg)


\le b\sup\limits_{n\ge 1}p^{(2)}


\bigg(\sum_{k=1}^{n}((L_1^M y)_R^{-1})_k Mx_k\bigg).


\end{multline*}


Так как $(L_1^M)_R^{-1}$ --- линейный непрерывный оператор из $H_2$ в


$A_1(MX,H_2)$, то $\forall p^{(2)}\in P_2$ \; $\exists p^{(3)}\in P_2$,


$\exists b_1< +\infty: \forall d=(d_k)_k^{\infty}\in A_1(MX,H_2)$


$$


q^1_{p^{(2)}}(d)=


\sup\limits_{n\ge 1}p^{(2)}\bigg(\sum_{k=1}^nd_kMx_k\bigg)\le b_1p^{(3)}


\bigg(\sum_{k=1}^{\infty}d_kMx_k\bigg).


$$


Поэтому


$$


p^{(1)}(By)\le bb_1 p^{(3)}\bigg(\sum_{k=1}^{\infty}


((L_1^M y)_R^{-1})_k Mx_k\bigg)=bb_1p^{(3)}(y),


$$


что и доказывает непрерывность оператора $B$.$\quad\square$








\begin{cornonum}


Пусть последовательность $X$ \ $M$-инвариантна относительно пары


$(H_1,H_2)$, а пространства $A_1(MX_1,H_2)$ и $A_1(X,H_1)$ изоморфны


(топологически).


Для того чтобы последовательность $MX$ была ППС в $H_2$, необходимо, а в


случае,


когда $X$ --- ППС в $H_1$, и достаточно, чтобы $M$ был эпиморфизмом $H_1$


на $H_2$, имеющим ЛНПО.


\end{cornonum}





Анализируя доказательство предложения 5, нетрудно убедиться в том, что


если $MX$ --- ЭППС в $H_1$, то оператор $B=M_R^{-1}$ можно определить


конструктивно. Поэтому справедливо








\begin{prop}


Пусть имеют место исходные предположения предложения $5$ с той


лишь разницей, что $MX$ --- не только ППС, но и ЭППС в $H_2$. Тогда $M$ ---


эпиморфизм $H_1$ на $H_2$ с конструктивно определяемым ЛНПО.


\end{prop}








\begin{cornonum}


Пусть последовательность $X$ \ $M$-инвариантна относительно пары


$(H_1,H_2)$ и пространства $A_1(MX,H_2)$, $A_1(X,H_1)$ топологически


изоморфны.


Для того чтобы $MX$ была ЭППC в $H_2$, необходимо, а в случае, когда $X$


--- ЭППС


в $H_1$, и достаточно, чтобы $M$ был эпиморфизмом $H_1$ на $H_2$ с


конструктивно определяемым ЛНПО.


\end{cornonum}





Аналогичные результаты справедливы и для правильных АПС.








\begin{prop}


Пусть последовательность $X$ \ $M$-абсолютно инвариантна относительно


$(H_1,H_2)$ и $A_2(MX,H_2)$ топологически изоморфно $A_2(X,H_1)$. Пусть,


далее, МХ --- ПАПС в $H_2$. Тогда $M$ --- эпиморфизм $H_1$ на $H_2$, имеющий


ЛНПО.


\end{prop}





Это предложение доказывается по тому же плану, что и предложение 5.


Именно, т.\,к. $MX$ --- ПАПС в $H_2$,


то существует ЛНПО $(L_0^M)_R^{-1}$ для оператора представления


$$


L_0^M: A_2(MX,H_2)\to H_2\mid \forall d\in A_2(MX,H_2)\to


L_0^M(d)=\sum_{k=1}^{\infty}d_kMx_k\in H_2.


$$


При этом


$\forall y\in H_2$ \; $y=\sum\limits_{k=1}^{\infty}((L_0^My)_R^{-1})_kMx_k$


и ряд абсолютно сходится в $H_2$. Следовательно, 


ряд $\sum\limits_{k=1}^{\infty}((L_0^My)_R^{-1})_kx_k$


сходится абсолютно в $H_1$ и если $v$ --- его сумма, то $Mv=y$. При этом


$v=By$, где $B$ --- линейный оператор из $H_2$ в $H_1:


\forall y\in H_2$ \ $By=\sum\limits_{k=1}^{\infty}((L_0^My)_R^{-1})_kx_k$.


Для любой преднормы


$p^{(1)}\in P_1$ найдутся $b<\infty$ и $p^{(2)}\in P_2$ такие, что


$\forall d=(d_k)_{k=1}^{\infty}\in A_2(MX,H_2)$ \


$\sum\limits_{k=1}^{\infty}|d_k|p^{(1)}(x_k)\le b\sum\limits_{k=1}^


{\infty}|d_k|p^{(2)}(Mx_k)$. Отсюда


$$


p^{(1)}(v)=p^{(1)}\bigg(\sum_{k=1}^{\infty}((L_0^My)_R^{-1})_kx_k\bigg)\le \\


\sum_{k=1}^{\infty}|((L_0^My)_R^{-1})_k|p^{(1)}(x_k)\le


b\sum_{k=1}^{\infty}|((L_0^My)_R^{-1})_k|p^{(2)}(Mx_k).


$$


В силу непрерывности оператора


$(L_0^M)_R^{-1}$ \; $\forall p^{(2)}\in P_2$ \; $\exists p^{(3)} \in P_2$,


$\exists b_1<\infty: \forall h=(h_k)_{k=1}^{\infty}\in A_2(MX,H_2)$ \


$q^2_{p^{(2)}}(h):=\sum\limits_{k=1}^{\infty}|h_k|p^{(2)}(Mx_k)\le b_1p^{(3)}


\Big(\sum\limits_{k=1}^{\infty}h_k Mx_k\Big)$.


Поэтому


$$


p^{(1)}(By)\le bb_1p^{(3)}


\Big(\sum\limits_{k=1}^{\infty}((L_0^M y)_R^{-1})_k Mx_k\Big)=bb_1 p^{(3)}(y),


$$


и непрерывность оператора $B$ доказана.








\begin{cornonum}


Пусть последовательность $X$ \ $M$-абсолютно ивариантна относительно


пары $(H_1,H_2)$, а пространства $A_2(MX,H_2)$ и $A_2(X,H_1)$


топологически изоморфны.





Для того чтобы последовательность $MX$ была ПАПС в $H_2$, необходимо, а в


случае $X$ --- ПАПС в $H_1$, и достаточно, чтобы $M$ был эпиморфизмом


$H_1$ на $H_2$, имеющим ЛНПО.


\end{cornonum}





Справедлив и аналог предложения 6.








\begin{prop}


Пусть имеют место исходные предположения предложения $7$, но


$MX$ --- не только ПАПС, но и ЭПАПС в $H_2$. Тогда $M$ --- эпиморфизм


$H_1$ на $H_2$ с конструктивно определяемым ЛНПО.


\end{prop}








\begin{cornonum}


Пусть последовательность $X$ \ $M$-абсолютно инвариантна относительно


пары $(H_1,H_2)$, а пространства $A_2(MX,H_2)$ и $A_2(X,H_1)$ топологически


изоморфны. Для того чтобы $MX$ была ЭАПС в $H_2$, необходимо, а в случае,


когда


$X$ --- ЭПАПС в $H_1$, и достаточно, чтобы $M$ был эпиморфизмом $H_1$ на


$H_2$ с конструктивно определяемым ЛНПО.


\end{cornonum}








3. Изложенные в этом параграфе результаты показывают, что исходные


предположения


в предложениях 1--3 из ([3], \S\,2) в известной степени существенны для


справедливости их заключений. Отметим еще, что понятие $M$-инвариантности


и $M$-абсолютной инвариантности является естественным обобщением понятий


$\Pi$-инвариантности и $\Pi$-абсолютной инвариантности, введенных в


([3], \S\,4) в частной ситуации, когда $M$ является оператором сужения


$\Pi$ функций, определенных на множестве $Q_1\subset R^p$, $p\ge 1$, и


принадлежащих некоторому пространству


$H_1$, на множество $Q_2\subset Q_1$; при этом предполагается, что сужение


$\Pi$ действует непрерывно из $H_1$ в некоторое пространство $H_2$


функций, определенных на множестве $Q_2$.








\section{Представляющие системы экспонент с мнимыми


показателями\protect\\ в пространствах


бесконечно дифференцируемых функций}





1. Пусть $H_j$ --- ПОЛВП ($j=1,2$), $x_k\in H_1 \cap H_2$


$\forall k\ge 1$,


и $X$ --- ПС (или АПС) в $H_1$. Если $X$ будет ПС (соответственно АПС)


в $H_2$, то будем говорить, что ПС (соответственно АПС)


$X$-продолжима из $H_1$ в $H_2$.





Это определение обобщает понятие продолжимой ПС или АПС, введенное в ([2],


гл.\,II, \S\,2)


в более частной ситуации, когда $H_1\subset H_2$, а также понятия


внутрь-продолжимой ПС и АПС (см. там же, а также [7]).





Пусть $G$ --- произвольное открытое множество в $\Bbb{R}^p$ и


$BC^\infty(G)$ ---


пространство всех комплекснозначных бесконечно дифференцируемых в $G$ функций,


равномерно ограниченных и равномерно непрерывных в $G$ вместе с любой


своей частной


производной. Топология в $BC^\infty(G)$ задается счетным набором норм


$$


\|y\|_n=\sup\{|y^{(\alpha)}(X)|:X\in G,\ |\alpha|_p\le n\},\ \


n=0,1,2,\dotsc,


$$


где $\alpha=(\alpha_1,\dotsc,\alpha_p)\in N_0^p$, $N_0=\{0,1,\dotsc\}$,


$|\alpha|_p=\alpha_1+\alpha_2+\dotsb+\alpha_p$.





В этой топологии $BC^\infty(G)$ --- пространство Фреше.


\medskip








2. Сформулируем относительно сходимости ряда по


общей системе экспонент с мнимыми показателями


\begin{equation}


\sum_{|k|_p=0}^\infty


c_k\exp\bigg(i\sum_{j=1}^{p}\mu_{k,j}x_j \bigg),\ \


\mu_{k,j}\in\Bbb{R},\ \ 1\le j\le p,\ \ k=(k_1,\dotsc, k_p)\in Z^p,


\end{equation}


где $Z:=(0,\pm 1,\pm 2,\dotsc)$, $|k|_p=|k_1|+\dotsb+|k_p|$, следующие


уиверждения:





1) pяд (2) сходится абсолютно в $BC^\infty(G)$ для любого открытого


множества $G$ в $\Bbb{R}^p$;





2) pяд (2) сходится абсолютно в $BC^\infty(G_1)$ для некоторого открытого


множества $G_1$ в $\Bbb{R}^p$;





3) pяд (2) сходится абсолютно в $BC^\infty(\Bbb{R}^p)$;





4) $\forall m\ge0$ \; $\sum\limits_{|k|_p=0}^\infty|c_k|\cdot|\mu_k|^m<\infty$,


где


$\mu_k=(\mu_{k,1},\dotsc,\mu_{k,p})$,


$|\mu_k|^m=|\mu_{k,1}|^m\cdot\dotsc\cdot|\mu_{k,p}|^m$.








\begin{prop}


Утверждения $1)$--\,$4)$ равносильны.


\end{prop}





Это предложение доказывается так же, как лемма 2.1 из [8], и основную роль


в доказательстве играет равенство


$|\exp i\langle\mu_k,X\rangle|=1$ \; $\forall X\in\Bbb{R}^p$,


$\forall \mu_k\in\Bbb{R}^p$


(здесь $\langle\mu_k,X\rangle =\sum\limits_{j=1}^p\mu_{k,j}x_j$).





Пусть теперь $T_{a,b}$ --- какой-либо прямоугольный открытый параллелепипед


$T_{a,b}=\{X=(x_1,\dotsc,x_p):a_k<x_k<b_k$, $k=1,\dotsc,p\}$,


$-\infty<a_k<b_k<+\infty$, $k=1,\dotsc,p$.





Рассмотрим специальную систему экспонент с мнимыми показателями


$$


E^T_p=\bigg\{\exp 2\pi i\sum_{j=1}^p\frac{k_jx_j}{b_j-a_j},


\ k_s=0,\pm 1,\pm 2,\dotsc, \ s=1,2,\dots,p\bigg\}=:


\big\{\exp\big\langle


2\pi ik,\ \tfrac{X}{b-a}\big\rangle\big\}^{\infty}_{|k|_p=0}.


$$





Из предложения 9 следует, что система $E^T_p$ \ $\Pi$-абсолютно инвариантна


относительно пары


$BC^\infty(T_{a,b})$, $BC^\infty(G)$, где $\Pi$ --- оператор сужения


функции из


$BC^\infty(T_{a,b})$ на $G$, если $G\subseteq T_{a,b}$. Нетрудно также


показать, пользуясь леммой 2.1 из [8], что


$E_p^T$ \ $\Pi$-инвариантна относительно той же пары.





Обозначим символом $B_T$ множество всех пар $(X,Y)$ различных точек


$X(x_1,\dotsc,x_p)$, $Y(y_1,\dotsc, y_p)$ из $\overline{T}_{a,b}$ таких,


что  $y_k-x_k=0,\pm (b_k-a_k)$ при $k\le p$. Положим еще


$$


H_0^T:=\{y(X)\in


BC^\infty(T_{a,b}):y^{(\alpha)}(X)=y^{(\alpha)}(Y) \ \forall \alpha\in


N_0^p, \ \forall (X,Y)\in B_T\}.


$$





В индуцированной из $BC^\infty(T_{a,b})$ топологии $H^T_0$ --- его


замкнутое пространство и, следовательно, тоже пространство Фреше.





Как нетрудно проверить, для каждой функции $y(X)$ из $H_0^T$ найдется


единственное сходящееся в


$BC^\infty(T_{a,b})$ разложение по системе $E_p^T$:


\begin{equation}


y(X)=\sum_{|k|_p=0}^\infty y_k\exp\big\langle 2\pi ik,\ \tfrac{X}{b-


a}\big\rangle.


\end{equation}


При этом (3) --- это ряд Фурье функции $y(X)$ по системе $E^T_p$. Он


сходится абсолютно в


$BC^\infty(T_{a,b})$, а его коэффициенты $y_k$ определяются эффективно по


известным формулам.


Все эти факты легко установить, применив, как в ([9], с.\,506), обычный


прием интегрирования


по частям интегрального представления для $y_k$ и использовав обращение в нуль


всех получающихся при этом


внеинтегральных членов для любой функции из $H^T_0$. Таким образом,


$E^T_p$ --- эффективный


абсолютный базис Шаудера в $H^T_0$. Возникает вопрос о том,


будет ли $E^T_p$


представляющей системой в $BC^\infty(G)$ и какого типа (ЭПС, ППС, ЭППС,


АПС и т.\,д.). Здесь и далее


$G$ --- открытое ограниченное множество в $\Bbb{R}^p$.


\medskip





3. Укажем вначале одно необходимое условие продолжимости $E^T_p$


из $H^T_0$ в $BC^\infty(G)$. Пусть $C(\overline{G})$ --- банахово пространство


непрерывных на компакте $\overline{G}$ функций с нормой


$\|y\|_0=\max\{|y(X)|:X\in\overline{G}\}$.








\begin{thm}


Если система $E^T_p$ полна в $C(\overline{G})$, то


$(\overline{G}\times \overline{G})\bigcap B_T=\emptyset$.


\end{thm}








{\bf Доказательство.} Допустим, что $(X_0,


Y_0)\in(\overline{G}\times\overline{G})\bigcap B_T$.


Так как $X_0=(x_{1,0},\dotsc, x_{p,0})\ne Y_0=(y_{1,0},\dotsc,y_{p,0})$,


то $\exists k_0\le p:


y_{k_0,0}-x_{k_0,0}=\pm(b_{k_0}-a_{k_0})$. Пусть для определенности


$y_{k_0,0}-x_{k_0,0}=


b_{k_0}-a_{k_0}$. Рассмотрим функцию $y_1(X)=x_{k_0}$ из $BC^\infty(G)$.


Ясно, что $y_1(X_0)\ne y_1(Y_0)$. В то же время в силу полноты $E^T_p$ в


$C(\overline{G})$ найдется


последовательность функций $\{g_n(X)\}_{n=1}^\infty$:


$$


g_n(X)=\sum_{k_j=-m_n}^{+m_n}a_{k,n}\exp\big\langle2\pi


ki,\ \tfrac{X}{b-a}\big\rangle,\ \ n=1,2,\dotsc;


\ \ m_n\in N_0,\ \ k=(k_1,\dots,k_p),


$$


сходящаяся к $y_1(X)$ равномерно на $\overline{G}$. Отсюда


$$


y_1(X_0)=\lim_{n\to\infty}g_n(X_0)=\lim_{n\to\infty}g_n(Y_0)=y_1(Y_0),


$$


что невозможно.$\quad\square$    %% put {pf}





Так как $y_1(X)\in BC^\infty(G)$ и $BC^\infty(G)\hookrightarrow


C(\overline{G})$, то точно так же


получаем








\begin{thm}


Если система $E^T_p$ полна в $BC^\infty(G)$, то


$(\overline{G}\times\overline{G})\bigcap B_T=\emptyset$.


\end{thm}








\begin{cornonum}


Если $G\subseteq T$ и $E^T_p$ --- ПС в $BC^\infty(G)$, то


$(\overline{G}\times\overline{G})\bigcap B_T=\emptyset$.


\end{cornonum}








\begin{thm}


Если $G\subseteq T$ и $(\overline{G}\times\overline{G})


\bigcap B_T=\emptyset$, то


\begin{itemize}


\item[1)] $E^T_p$ --- ПС в $BC^\infty(G)$ тогда и только тогда,


когда $\mathrm a_1)$~существует


непрерывный оператор продолжения любой функции из $BC^\infty(G)$ до


функции


из $H_0^T$ $($т.\,е. непрерывный оператор 


продолжения $BC^\infty(G)$ в $H_0^T);$


\item[2)] $E^T_p$ --- ЭПС в $BC^\infty(G)$ тогда и только тогда, 


когда $\mathrm a_2)$~можно


конструктивно определить оператор продолжения $BC^\infty(G)$ в $H_0^T;$


\item[3)] $E^T_p$ --- ППС в $BC^\infty(G)$ тогда и только тогда, 


когда $\mathrm a_3)$~существует линейный непрерывный оператор 


продолжения $BC^\infty(G)$ в $H_0^T;$


\item[4)] $E^T_p$ --- ЭППС в $BC^\infty(G) \iff \mathrm a_4)$ можно определить


конструктивно


линейный непрерывный оператор продолжения $BC^\infty(G)$ в $H_0^T;$


\item[5)] $E^T_p$ --- АПС в $BC^\infty(G) \iff$ выполнено $\mathrm a_1);$


\item[6)] $E^T_p$ --- ЭАПС в $BC^\infty(G) \iff$ имеет место $\mathrm a_2);$


\pagebreak





\item[7)] $E^T_p$ --- ПАПС в $BC^\infty(G) \iff$ справедливо $\mathrm a_3);$


\item[8)] $E^T_p$ --- ЭПАПС в $BC^\infty(G) \iff$ выполнено $\mathrm a_4)$.


\end{itemize}


\end{thm}








{\bf Доказательство.} Так как топология в $H_0^T$ индуцируется


из $BC^\infty(T_{a,b})$, то $E^T_p$ \ $\Pi$-инвариантна и


$\Pi$-абсолютно инвариантна


относительно пары $H_0^T$, $BC^\infty(G)$. Кроме того, если оператор


сужения $\Pi$ пространства $H_0^T$ на $BC^\infty(G)$ является эпиморфизмом


и если $B$ --- какой-либо его правый обратный оператор, то $\forall y\in


BC^\infty(G)$ \; $\Pi By=By |_G=y$. Таким образом, любой такой оператор


$B$ является оператором продолжения $BC^\infty(G)$ в $H_0^T$. При этом,


т.\,к. $H^T_0$, $BC^\infty(G)$ --- пара Майкла, то $\Pi$ отображает $H^T_0$ на


$BC^\infty(G)$


тогда и только тогда, когда существует непрерывный оператор продолжения


$BC^\infty(G)$ в $H^T_0$. Учитывая сказанное, заключаем, что


утверждения 1)--2) непосредственно вытекают из следствия~1 предложения~4;


утверждение~3) --- из следствия предложения~5; утверждение~4) --- из


следствия предложения~6; утверждения~5)--6) --- из следствия~2 предложения~4;


утверждение~7) --- из следствия предложения~7 и, наконец,


утверждение~8) --- из следствия предложения~8.








Для пространства $BC^\infty(\Bbb{R}^p)$ введем подпространство


$$


H_0^{\Bbb{R}^p}:=\{f(X)\in


BC^\infty(\Bbb{R}^p):f^{(\alpha)}(X)=f^{(\alpha)}(Y) 


\ \forall \alpha\in N_0^P,\ \forall(X,Y)\in B_{\Bbb{R}^p}\},


$$


где $B_{\Bbb{R}^p}$ --- множество всех пар различных точек $X$, $Y$ из


$\Bbb{R}^p$ таких,


что  $y_k-x_k=l_k(b_k-a_k)$ при $k=1,2,\dotsc,p$, где $l_k\in


Z=\{0,\pm 1,\pm 2,\dotsc\}$. Подпространство $H_0^{\Bbb{R}^p}$ в топологии, 


индуцированной из $BC^\infty(\Bbb{R}^p)$, является пространством


Фреше. Положим $H_1=H_0^{\Bbb{R}^p}$, $H_2=BC^\infty(G)$, где $G$ ---


открытое множество в


$\Bbb{R}^p$. Тогда $H_1,\ H_2$ --- пара Майкла,


а $E^T_p$ \ $\Pi$-инвариантна и $\Pi$-абсолютно


инвариантна относительно этой пары. Здесь $\Pi$ --- оператор сужения


значений любой функции из $H_0^{\Bbb{R}^p}$ на множество $G$.





Точно такими же рассуждениями, как в случае теоремы 3, и при тех же


исходных предположениях


устанавливаем полный ее аналог --- теорему~$3'$, для формулировки которой


достаточно лишь заменить $H_0^T$ на $H_0^{\Bbb{R}^p}$ в теореме 3.


\medskip





4. Пусть $G$ --- открытое множество, лежащее внутри $T_{a,b}$, и


пусть $d=\rho(


\overline{G},\partial T_{a,b})$. Допустим, что имеется оператор $M_0$


продолжения из $BC^\infty


(G)$ в $H_0^T$. Как известно (см., напр., [10], гл.\,I, \S\,1, п.\,1.4),


можно (эффективно) построить


функцию $b(t)$ из подпространства $C^\infty_0(\Bbb{R}^p)$ с компактным


носителем такую, что $b|_{\overline{G}}=1$ и носитель функции $b(X)$ 


находится в $(G)_{d/2}:=\{X\in


\Bbb{R}^p:\rho(X,\overline{G})\le d/2\}$. Тогда оператор


$B_0:B_0y=b(X)M_0y$ отображает


$BC^\infty(G)$ в $C^\infty_0(T_{a,b})\subset H_0^T$. Нетрудно заметить,


что если $M_0$


(как оператор из $BC^\infty(G)$ в $H_0^T$) непрерывен, линеен или


конструктивно определен, то


тем же свойством обладает и $B_0$ (как оператор из $BC^\infty(T_{a,b})$ и


из $BC^\infty(G)$


в $H_0^T$). Кроме того, оператор $B_0$ отображает $BC^\infty(G)$ в


$C_0^\infty(\Bbb{R}^p)$


(и в $BC^\infty(\Bbb{R}^p)$) и линеен, непрерывен или


конструктивно определен, если соответствующим свойством обладает 


оператор $M_0$: $BC^\infty(G)\to H^T_0$.





Пусть, обратно, имеется оператор $M_1$ продолжения из $BC^\infty(G)$ в


$BC^\infty(\Bbb{R}^p)$. Тогда $B_1y:=b(X)M_1y$ --- оператор продолжения


из $BC^\infty(G)$ в $H_0^T$ (и из $BC^\infty(G)$ в


$C_0^\infty(\Bbb{R}^p)$), который


обладает свойством непрерывности, линейности, конструктивной


определимости, если


соответствующее свойство имеется у $M_1$. Используя теорему~3 и ее


аналог (теорему $3'$) для $H_0^{\Bbb{R}^p}$, а также $\Pi$-инвариантность


и $\Pi$-абсолютную инвариантность системы $E^T_p$ относительно всех пар


($BC^\infty(G)$,$H_0^T$), ($BC^\infty(G)$,$H_0^{\Bbb{R}^p}$),


($BC^\infty(G)$,


$C_0^\infty(T_{a,b})$), ($BC^\infty(G)$,$BC^\infty(\Bbb{R}^p)$),


приходим к такому результату.








\begin{thm}


Пусть $G$ --- открытое множество, лежащее внутри $T_{a,b}$.


Тогда в каждой из приведенных ниже четырех групп утверждений 


все они равносильны друг другу.





$\mathrm{I}.1)$ $E^T_p$ --- ПС в $BC^\infty(G);$





$\mathrm{I}.2)$ $E^T_p$ --- АПС в $BC^\infty(G);$





$\mathrm{I}.3)$ существует оператор продолжения из $BC^\infty(G)$ в $H_0^T;$





$\mathrm{I}.4)$ имеется оператор продолжения из $BC^{\infty}(G)$ в


$C_0^\infty(T_{a,b});$





$\mathrm{I}.5)$ существует оператор продолжения из $BC^\infty(G)$ в $C_0^{\infty}


(\Bbb{R}^p);$





$\mathrm{I}.6)$ существует оператор продолжения из $BC^{\infty}(G)$ в


$H_0^{\Bbb{R}^p};$





$\mathrm{I}.7)$ имеется оператор продолжения из $BC^\infty(G)$ в


$BC^\infty(\Bbb{R}^p)$.


\smallskip





$\mathrm{II}.1)$ $E^T_p$ --- ЭПС в $BC^\infty(G);$





$\mathrm{II}.2)$ $E^T_p$ --- ЭАПС в $BC^\infty(G);$





$\mathrm{II}.3)$--$\mathrm{II}.7)$ оператор


из $\mathrm{I}.3)$--$\mathrm{I}.7)$ можно определить конструктивно.


\smallskip





$\mathrm{III}.1)$ $E^T_p$ --- ППС в $BC^\infty(G);$





$\mathrm{III}.2)$ $E^T_p$ --- ПАПС в $BC^\infty(G);$





$\mathrm{III}.3)$--$\mathrm{III}.7)$ оператор продолжения 


из $\mathrm{I}.3)$--$\mathrm{I}.7)$ линеен и непрерывен.


\smallskip





$\mathrm{IV}.1)$ $E^T_p$ --- ЭППС в $BC^\infty(G);$





$\mathrm{IV}.2)$ $E^T_p$ --- ЭПАПС в $BC^\infty(G);$





$\mathrm{IV}.3)$--$\mathrm{IV}.7)$ оператор


из $\mathrm{I}.3)$--$\mathrm{I}.7)$


линеен, непрерывен и может быть конструктивно определен.


\end{thm}








\begin{notenonum}


Из теоремы 2.4 работы [8] следует, что любое из утверждений


I.1)--I.7) равносильно такому: I.8)~в $BC^\infty(G)$ имеется хотя бы одна АПС


экспонент с чисто мнимыми показателями


$E_{\mu}=\{\exp\langle i\mu_k,X\rangle\}^{\infty}_{|k|_p=0}$,


$\mu_k\in \Bbb{R}^p$.


\end{notenonum}





5. Пусть $X_0\in\Bbb{R}^p$ и


$C^\infty_{X_0}:=\ind\limits_


{\overrightarrow{\ m\ }} C^\infty(Q_m)$, $Q_m=\{X\in\Bbb{R}^p:|X-X_0|_p


\leqslant\frac{1}{m}\}$, $m=1,2,\dotsc$ Пусть, далее,


$E_{\mu}:= \{\exp\langle i\mu_k,X\rangle 


\}^\infty_{|k|_p=0}$, где $\mu_k\in\Bbb{R}^p$ \; $\forall k\ge 1$.





Рассмотрим какое-либо полное отделимое бочечное пространство $H$


бесконечно дифференцируемых функций такое, что


\begin{equation}


BC^{\infty}(\Bbb{R}^p)\hookrightarrow H\hookrightarrow


C^\infty_{X_0}.


\end{equation}





Из леммы 2.1 работы [8] следуют равенства


\begin{multline*}


A_2(E_{\mu},C^{\infty}_{X_0})=A_2(E_{\mu},


BC^{\infty}(\Bbb{R}^p))=\bigg\{d=(d_k)^{\infty}_{|k|_p=0}: \\


%


\max \bigg[\sum^\infty_{|k|_p=0}|d_k|\,|(\mu_k)^{\alpha}|:


|\alpha|_p\le m,\ \alpha\in N_0^p\bigg]=:[d]_m<\infty,


\ m=1,2,\dotsc\bigg\}.


\end{multline*}


Отсюда $A_2(E_{\mu},C_{X_0}^{\infty})=A_2(E_{


\mu},H)=A_2(E_{\mu},BC^{\infty}(\Bbb{R}^p))$. Топология в каждом из этих трех


пространств, совпадающих по набору элементов, задается нормами $[d]_m$, и


все они --- пространства Фреше.





Предположим, что $E_{\mu}$ --- АПС в $H$. Тогда оператор представления


$L_0: \forall\alpha\in A_2$ \; $(E_{\mu},H)\to


\sum\limits_{\|k|_p=0}^{\infty}d_k\exp\langle i\mu_{k},X\rangle \in H$ 


является эпиморфизмом $A_2(E_{\mu},H)$ на $H$. Так как


$A_2(E_{\mu},H)$ совершенно полно, а $H$ бочечно, то по теореме об открытом


отображении (см. [11], с.\,170, теорема~7) $L_0$ открыто. Но тогда


по предложению 13 из ([11], с.\,174) $H$ --- пространство Фреше. Таким


образом, справедлива








\begin{thm}


Пусть $H$ --- полное отделимое ЛВП со свойством $(4)$ и пусть в $H$


имеется хотя бы одна АПС экспонент c чисто мнимыми показателями. Тогда $H$


--- пространство Фреше.


\end{thm}





Напомним еще, что теорема 5 не допускает прямого обращения.


Действительно, если $G$ --- открытое множество в $\Bbb{R}^p$,


$G=\bigcup\limits_{n=1}^{\infty}Q_n$, где $\forall n


\ge 1$ \; $Q_n$ --- толстый компакт (определение толстого компакта можно


найти,


напр., в [8], [9], [12]) и $Q_n\subset Q_{n+1}\subset G$, то, как


показано в [8], в пространстве Фреше $C^\infty(G)=\proj\limits_{


\overleftarrow{\ n\ }}C^{\infty}(Q_n)$ 


нет ни одной АПС экспонент с чисто


мнимыми показателями. В связи с этим возникают задачи о характеризации


пространства Фреше, имеющего хотя бы одну АПС экспонент вида $E_{\mu}$, а


также установления критерия того, что какая-либо фиксированная система


$E_{\mu}$ является АПС в данном пространстве Фреше бесконечно-дифференцируемых


функций. Возможно, что эти задачи можно решить с помощью методов и результатов


статей [13], [14], но это уже выходит за рамки данной работы.











\begin{thebibliography}{99}





\bibitem{1}


Коробейник Ю.Ф. {\em Об одной двойственной задаче}. I. {\em Общие результаты.


Приложения к пространствам Фреше} // Матем. сб. -- 1975. -- Т.\,97.


-- \No\,2. -- C.\,193--226.


\bibitem{2}


Коробейник Ю.Ф. {\em Представляющие системы} // УМН. -- 1981.


-- Т.\,36. -- Вып.\,2. -- С.\,73--126.


\bibitem{3}


Коробейник Ю.Ф. {\em О некоторых классах представляющих систем и их


преобразованиях}. I // Тр. Матем. центра им.~Н.И.\,Лобачевского. --


Казань, Казанское матем. об-во. -- 2002. -- Т.\,14. -- С.\,171--183.


\bibitem{4}


Канторович Л.В., Акилов Г.П. {\em Функциональный анализ в нормированных


пространствах}. -- М.: Физматгиз, 1959. -- 684~с.


\bibitem{5}


Епифанов О.В. {\em О существовании непрерывного правого обратного оператора


в одном классе локально выпуклых пространств} // Изв. Северо-Кавказск.


Научного центра высшей школы. Сер.


естеств. наук. -- 1991. -- \No\,3. -- C.\,3--4.


\bibitem{6}


Michael E. {\em Continuous selections}. I // Annals of Math.


-- 1956. -- V.\,63. -- \No\,2. -- P.\,361--382.


\bibitem{7}


Коробейник Ю.Ф., Леонтьев А.Ф. {\em О свойствах внутрь-продолжаемости систем


экспонент} // Матем. заметки. -- 1980. -- Т.\,28. -- Вып.\,2. -- С.\,243--254.


\bibitem{8}


Korobeinik Yu.F. {\em On absolutely representing systems in spaces of


infinitely differentiable functions} // Studia


Math. -- 2000. -- V.\,139. -- \No\,2. -- P.\,175--188.


\bibitem{9}


Korobeinik Yu.F. {\em Representing systems of exponentials in the spaces of


infinitely-differentiable functions and extendability in the sense of


Whitney} //


Тurkish Journ. of Math. -- 2001. -- V.\,25. -- \No\,4. -- P.\,503--517.


\bibitem{10}


Х\"ермандер Л. {\em Анализ линейных дифференциальных операторов с частными


производными}. Т.\,1. {\em Теория распределений и анализ Фурье}.


-- М.: Мир, 1986. -- 462~с.


\bibitem{11}


Робертсон А.П., Робертсон В.Дж. {\em Топологические векторные 


пространства}. -- М.: Мир, 1967. -- 256~с.


\bibitem{12}


Коробейник Ю.Ф. {\em О толстых компактах в $\Bbb{R}^p$} // Изв.


вузов. Северо-Кавказский регион. -- 2002. -- Юбилейный выпуск. -- С.\,93--94.


\bibitem{13}


Коробейник Ю.Ф. {\em Абсолютно представляющие семейства} // Матем. заметки.


-- 1987. -- Т.\,42. -- \No\,5. -- С.\,670--680.


\bibitem{14}


Коробейник Ю.Ф. {\em Абсолютно представляющие семейства и реализация


сопряженного пространства} // Изв. вузов. Математика.


-- 1990. -- \No\,2. -- C.\,68--76.





\end{thebibliography}





\vskip 2cm





\post{\it Ростовский государственный}{\it Поступила}


\post{\it университет}{01.07.2003}





\label{end}


\end{document}





